\documentclass[../main.tex]{subfiles}
\begin{document}

\chapter{Conclusion}
\label{ch:conclusion}

\section{Summary of Contributions}

This thesis investigated whether multi-agent LLM architectures can provide intrusion detection that balances operational accuracy with model transparency. The experimental results demonstrate that the proposed system satisfies both objectives.

The first contribution is the first systematic, per-attack-type evaluation of a multi-agent LLM intrusion detection system at realistic class distributions. Prior evaluations of LLM-based NIDS used balanced datasets or aggregated multi-attack batches, neither of which reflects the class imbalance of production networks. The Stage 1 evaluation reveals a performance range from F1=100\% on brute-force and DoS categories to F1=0\% on Infiltration, a range that would be invisible in aggregate evaluation. The necessity of specialist agents is confirmed by an ablation study demonstrating that the Temporal Agent alone accounts for a 44 percentage point recall difference on SSH-Bruteforce and 36 percentage points on DDoS-HOIC, and that a two-agent configuration degrades to 0--10\% recall across all tested attack types. The HOIC extension reveals that the fixed 30\% weight assigned to the Devil's Advocate agent suppresses correct detections on ambiguous attacks, improving recall by 40 percentage points when removed, demonstrating that optimal agent configuration is attack-type-dependent.

The second contribution is a two-tier ML-plus-LLM architecture reducing the cost of LLM-based detection by 94.6 per cent. Actual expenditure across all fourteen Stage 1 experiments was \$27.35, against an estimated \$509.78 without Tier 1 pre-filtering. The Random Forest at threshold 0.15 achieves one hundred per cent recall on attack flows in the development partition, demonstrating that operating costs can be reduced without degrading the detection rate. A random filter control experiment validates that the RF provides intelligent routing: a random 7\% selection achieves zero per cent recall whilst the trained RF at the same selection rate achieves 100 per cent, confirming that the cost reduction is inseparable from the routing intelligence. This two-tier pattern is generalisable to any domain where LLM analysis is accurate but expensive and a cheaper discriminator can identify the easy cases.

The third contribution is the introduction and empirical evaluation of a Devil's Advocate agent as a structural false positive control mechanism. The DA achieves zero per cent FPR across thirteen of fourteen attack categories, but the HOIC ablation reveals that the fixed 30\% weight is counterproductive on ambiguous attacks; removing the DA improves recall from 58\% to 98\% on DDoS-HOIC without increasing the false positive rate. This dual finding extends the multi-agent debate literature \autocite{Yin2024, Du2023, Chan2023} with the first evidence that adversarial agent effectiveness is non-monotonic in classification difficulty, and identifies confidence-gated adversarial weighting as a concrete engineering direction for multi-agent security systems.

The fourth contribution is a complete explainability framework where every detection decision is accompanied by multi-perspective reasoning chains from six analytical agents, operationalised through the Flow Inspector interface. Unlike SHAP or LIME post-hoc explanations, AMATAS reasoning chains are generated as part of the detection process itself, grounded in domain knowledge, and structured to present both evidence for malice and the best available counter-argument. A faithfulness audit verified that 89.8\% of factual claims in agent reasoning are correct against actual flow data, establishing that the reasoning chains are predominantly grounded in observed evidence rather than confabulated.

The fifth contribution is a controlled seven-configuration ablation study quantifying the marginal value of each architectural addition, from zero-shot single agent (F1=62.8\%) through prompt-engineered single agent (F1=58.0\%) and tool-augmented configurations (F1=36.7--58.3\%) to multi-agent AMATAS (F1=88\%). The extended comparison (Configurations D--G) demonstrates that even dataset-compatible tools that return genuinely informative results actively degrade single-agent performance, establishing that multi-agent debate rather than tool augmentation is the critical architectural contribution.

\section{Answers to the Research Questions}

\textbf{RQ1} asked whether a multi-agent LLM architecture can achieve detection accuracy comparable to traditional machine learning classifiers. AMATAS achieves mean F1 of 86.3\% across all fourteen attack types (92.9\% excluding Infiltration), which falls below the theoretical ceiling of fine-tuned discriminative models that report F1 exceeding 99\% on CICIDS2018. However, a standalone evaluation of the Tier~1 Random Forest (Section~\ref{sec:rf-baseline}) reveals that this comparison is misleading: the RF achieves mean F1 of only 57.1\%, with perfect detection on seven attack types but complete failure (0\% F1) on the remaining seven. AMATAS detects five of those seven RF-invisible categories at F1 $\geq$ 72\%, demonstrating that the multi-agent LLM pipeline provides detection capability on attack types where statistical classifiers fail entirely. The accuracy gap between AMATAS and fine-tuned models is the cost of generality and explainability: AMATAS produces multi-perspective reasoning chains for every verdict, a capability that no traditional classifier provides. The ablation study confirms that this performance requires multi-agent specialisation; a minimal two-agent configuration achieves only 0--12\% recall on brute force and 10\% on DDoS-HOIC, and the Temporal Agent's removal alone reduces recall by 36--44 percentage points. The HOIC ablation further reveals that the optimal agent configuration is attack-type-dependent: removing the Devil's Advocate improves HOIC recall from 58\% to 98\%, demonstrating that fixed adversarial weighting is suboptimal across the full attack taxonomy.

\textbf{RQ2} asked how a two-tier architecture affects the cost-accuracy trade-off. The Tier 1 Random Forest filters approximately ninety-four per cent of flows as confidently benign, reducing Stage 1 cost from an estimated \$509.78 to an actual \$27.35, a 94.6 per cent reduction, whilst maintaining one hundred per cent recall on attack flows at the operating threshold. The resulting cost of approximately \$0.002 per flow makes LLM-based analysis economically viable for batch processing of millions of flows.

\textbf{RQ3} asked about the quality and utility of multi-agent reasoning chains relative to opaque classifier outputs. The reasoning chains generated by AMATAS are qualitatively different from feature attribution outputs in several respects: they are generated as part of detection rather than as post-hoc approximations; they incorporate domain knowledge about attack signatures, protocol semantics, and MITRE ATT\&CK techniques; they present structured disagreement between analytical perspectives; and they include an explicit counter-argument. A faithfulness audit across 758 LLM-analysed flows verified that 89.8\% of 6{,}279 factual claims in agent reasoning were correct against actual flow data, with confabulation concentrated in TCP flag interpretation (22.4\% error rate) and protocol identification (19.4\%) rather than in the numeric features that dominate detection logic. These results demonstrate that AMATAS reasoning chains are predominantly grounded in observed evidence, and that the structural properties of the multi-agent framework (specificity, grounding, multi-perspectival structure, and falsifiability) are complemented by empirically verified factual accuracy.

\textbf{RQ4} asked how detection performance varies across attack types and what the fundamental limitations of flow-level LLM analysis are. Performance ranges from F1=100\% for attacks with distinctive flow signatures to F1=0\% for Infiltration, where individual attack flows are statistically indistinguishable from benign traffic. No single-flow analysis method, LLM or otherwise, can reliably detect attacks that have been designed to resemble benign traffic at the individual flow level. The v3 temporal clustering experiment (Section~\ref{sec:v3-clustering}) partially addresses this limitation, recovering 58\% recall on Infiltration by presenting agents with multi-flow clusters rather than isolated records, though at the cost of an elevated 19.2\% false positive rate that underscores the fundamental difficulty of this attack category.

\section{Future Work}
\label{sec:future}

The AMATAS v3 temporal clustering experiment (Section~\ref{sec:v3-clustering}) demonstrated that grouping flows by source IP recovers substantial detection capability on Infiltration, raising recall from 0\% to 58\% under the best condition. However, the elevated false positive rate of 19.2\% (nearly four times the worst-case FPR observed in the remaining thirteen Stage 1 experiments) indicates that the clustering mechanism requires further refinement. Promising directions include adaptive clustering thresholds that vary the DNS-query trigger based on historical IP behaviour, post-clustering confidence calibration to suppress false positives without sacrificing recall, and a dedicated Infiltration-specific agent trained on multi-stage intrusion patterns. Additionally, extending the clustering evaluation to DDoS-HOIC and other attack types where temporal context may improve detection would establish the generalisability of the v3 architecture beyond the single category tested here.

The ablation study identifies adaptive agent weighting as an immediate priority. The Devil's Advocate's fixed 30 per cent contribution suppresses recall on DDoS-HOIC by 40 percentage points relative to the no-DA configuration, whilst leaving brute force recall unchanged. A weighting scheme that scales agent contributions according to inter-agent confidence variance (reducing the DA's weight when specialist agents are in strong agreement and increasing it when they diverge) would address this limitation without sacrificing the false positive control that motivates the adversarial mechanism. This represents the most immediately actionable architectural improvement identified by the evaluation. Extending the ablation to the remaining eleven categories, particularly Bot (F1=85\%) and the DoS variants, would establish whether the DA's counterproductive effect on HOIC generalises to all moderate-confidence attack types, and the four-agent configuration's 90\% recall on HOIC at 49\% cost motivates investigation of attack-type-adaptive agent selection, where a meta-classifier selects the optimal agent subset based on the Tier~1 RF's confidence profile.

Evaluation across alternative frontier LLMs is a priority for external validity. Claude Sonnet 4, Llama-3.3-70B, and Gemini 1.5 Pro are candidates for comparative evaluation; each would require re-calibration of the Devil's Advocate weight and Orchestrator consensus thresholds against the model's specific reasoning and confidence patterns. \textcite{Ferrag2020} document substantial performance variation across deep learning architectures on the same dataset, and equivalent or greater sensitivity to model choice should be expected for LLM-based systems where reasoning patterns differ more fundamentally between model families.

A production deployment integrating MCP-based tool access against live traffic with real IP addresses represents a natural extension of the comparison experiments. The null result observed on anonymised data does not preclude substantial tool utility on production traffic, where AbuseIPDB reputation scores, geolocation data, and MITRE ATT\&CK technique lookups could provide contextual enrichment that the agents' parametric knowledge cannot. Such a deployment would require addressing the security considerations discussed in Section~\ref{sec:mcp-security}, particularly tool response verification and permission boundaries, however, the MCP server infrastructure developed for the comparison experiments can be adapted to support this extension.

The held-out test partition (8.05 million flows excluded from all development and validation decisions) serves as the primary evaluation of AMATAS generalisability and should be treated as a one-shot evaluation with no post-hoc parameter adjustment. Longer-term directions include LoRA-adapted specialist models replacing each agent with a fine-tuned variant trained on agent-specific analytical labels, a streaming deployment architecture for real-time flow processing with asynchronous LLM queuing, and a formal user study providing empirical grounding for the explainability claims that this thesis argues for on structural grounds.

\onlystandalone{\printbibliography}
\end{document}
